% =========================================================================
% Academic Mathematical Problem Sheet / Examination Paper
% Powered by Jinwen-XU/homework (CTAN / TeX Live / MiKTeX Standard Class)
% Clean, minimal, publication-grade academic layout (Zero SaaS UI clutter)
% =========================================================================

\documentclass[
  a4paper,
  11pt,
  title in boldface,
  theorem in new line,
  remove problem qed
]{homework}

% 设置语言环境为中文 (支持中英文混排、经典定理与题型名称本地化)
\UseLanguage{Chinese}

% 常用数学与排版增强宏包
\usepackage{amsmath,amssymb,mathtools,bm}
\usepackage{tasks}      % 专业选择题多列对齐宏包
\usepackage{booktabs}   % 经典学术三线表宏包
\usepackage{array}

% 公式字体：现代学术 Typst 同款字体 (New Computer Modern Math)
\setmathfont{NewCMMath-Book.otf}

% 极简纯净页码设置（单次编译即显示正确页码）
\usepackage{fancyhdr}
\pagestyle{fancy}
\fancyhf{}
\fancyfoot[C]{\small\normalfont 第 \thepage\ 页}
\renewcommand{\headrulewidth}{0pt}
\renewcommand{\footrulewidth}{0pt}

% 选择题 tasks 标签格式设置为 A. B. C. D.
\settasks{
  label = \Alph*.,
  label-width = 1.6em,
  label-format = {\bfseries},
  item-indent = 2.2em,
  before-skip = 0.3em,
  after-skip = 0.5em
}

% 文档元数据（标准学术卷头）
\title{工科数学分析 · 第 1 章 极限与连续 \\
  \large\normalfont 《大邮数学集》期中与期末真题精选习题册}
\author{\small\normalfont 许可协议：CC BY-NC-SA 4.0}
\date{\today}

\begin{document}

\maketitle
% =========================================================================
% 一、选择题
% =========================================================================
\section*{一、选择题（下列各题给出的四个选项中，只有一个选项符合题目要求。）}

\begin{problem}
  已知非负数列 $\{a_n\}, \{b_n\}, \{c_n\}$．且 $\lim_{n \to \infty} a_n = 0$, $\lim_{n \to \infty} b_n = 1$, $\lim_{n \to \infty} c_n = +\infty$，则 (\quad)．

  \begin{tasks}(2)
  \task $a_n < b_n$
  \task $b_n < c_n$
  \task $\lim_{n \to \infty} a_n c_n$ 不存在
  \task $\lim_{n \to \infty} b_n c_n$ 不存在
\end{tasks}
\end{problem}

\begin{problem}
  设 $f(x) = \frac{1}{x^2} \sin x$, $g(x) = \frac{1}{x}$，则当 $x \to \infty$ 时，$f(x)$ 是 $g(x)$ 的 (\quad)．

  \begin{tasks}(2)
  \task 高阶无穷小
  \task 低阶无穷小
  \task 等价无穷小
  \task 同阶但非等价无穷小
\end{tasks}
\end{problem}

\begin{problem}
  若 $\lim_{x \to a} f(x) = \infty$, $\lim_{x \to a} g(x) = \infty$，则必有 (\quad)．

  \begin{tasks}(2)
  \task $\lim_{x \to a} (f(x) + g(x)) = \infty$
  \task $\lim_{x \to a} (f(x) - g(x)) = 0$
  \task $\lim_{x \to a} \frac{f(x)}{g(x)} = \infty$
  \task $\lim_{x \to a} k f(x) = \infty, k \ne 0$
\end{tasks}
\end{problem}

\begin{problem}
  $\lim_{x \to 0^+} \left(\frac{1}{x}\right)^{\sin x} = (\quad)$．

  \begin{tasks}(4)
  \task 0
  \task 1
  \task $e$
  \task 不存在
\end{tasks}
\end{problem}

\begin{problem}
  设 $f(x)$ 在 $x=x_0$ 的某个邻域内有定义，下面的条件中为 $f(x)$ 在 $x=x_0$ 处可导的充分必要条件是 (\quad)．

  \begin{tasks}(1)
  \task $\lim_{h \to 0} \frac{f(x_0+h) - f(x_0)}{h}$ 存在
  \task $\lim_{h \to 0} \frac{f(x_0+h^2) - f(x_0)}{h^2}$ 存在
  \task $\lim_{\Delta x \to 0} \frac{f(x_0+\Delta x) - f(x_0-\Delta x)}{\Delta x}$ 存在
  \task $\lim_{\Delta x \to 0} \frac{f(x_0) - f(x_0-\Delta x)}{\Delta x}$ 存在
\end{tasks}
\end{problem}

\begin{problem}
  当 $x \to +\infty$ 时，下列函数中与 $\sqrt{x^2+1}-x$ 为等价无穷小的是 (\quad)．

  \begin{tasks}(4)
  \task $\sin \frac{1}{x}$
  \task $\ln (1-\frac{1}{x})$
  \task $1-\cos \frac{1}{x}$
  \task $e^{1/x^2}-1$
\end{tasks}
\end{problem}

\begin{problem}
  设有数列：\textcircled{\scriptsize 1} $\left\{ (2^n+(-2)^n)^{1/n} \right\}$，\textcircled{\scriptsize 2} $\left\{ \frac{1}{1+2} + \frac{1}{1+2^2} + \dots + \frac{1}{1+2^n} \right\}$，\textcircled{\scriptsize 3} $\left\{ (1+\frac{1}{n})^n \right\}$，\textcircled{\scriptsize 4} $\left\{ \frac{n+2^n}{3^n} \right\}$，其中收敛数列的个数为 (\quad)．

  \begin{tasks}(4)
  \task 1
  \task 2
  \task 3
  \task 4
\end{tasks}
\end{problem}

\begin{problem}
  设 $f(x) = \begin{cases} \frac{2+e^{1/x}}{1+e^{1/x}} + \arctan \frac{1}{x}, & x \ne 0 \\ 0, & x=0 \end{cases}$，则 $x=0$ 是 $f(x)$ 的 (\quad)．

  \begin{tasks}(2)
  \task 连续点
  \task 第一类间断点，且为跳跃间断点
  \task 第一类间断点，且为可去间断点
  \task 第二类间断点
\end{tasks}
\end{problem}

\begin{problem}
  设 $\lim_{x \to 0} \left( 1+\frac{f(2x)}{x} \right)^{1/x} = e$，则 $\lim_{x \to 0} \frac{f(x)}{x^2} = (\quad)$．

  \begin{tasks}(4)
  \task $\frac{1}{4}$
  \task $\frac{1}{2}$
  \task 2
  \task 4
\end{tasks}
\end{problem}

\begin{problem}
  设函数 $f(x) = \begin{cases} \frac{e^x-1-x}{x^2}, & x < 0 \\ \frac{1}{\ln(1+x)}, & x \ge 0 \end{cases}$，则 $f(x)$ 在 $x=0$ 处 (\quad)．

  \begin{tasks}(2)
  \task 不连续
  \task 连续，但不导
  \task 一阶可导，但二阶不可导
  \task 二阶可导
\end{tasks}
\end{problem}

% =========================================================================
% 二、填空题
% =========================================================================
\section*{二、填空题（把答案填在题中横线上。）}

\begin{problem}
  $\lim_{n \to \infty} (\sqrt{n+\sqrt{n}} - \sqrt{n-\sqrt{n}}) = \underline{\hspace{3.5em}}$．

  \vspace{1.2cm}
\end{problem}

\begin{problem}
  $\lim_{x \to \infty} \frac{3x^2+5}{5x+3} \sin \frac{2}{x} = \underline{\hspace{3.5em}}$．

  \vspace{1.2cm}
\end{problem}

\begin{problem}
  若 $\lim_{x \to 0} \frac{f'(x)}{g'(x)}$ 不存在，则由 LHôpital 法则得 $\lim_{x \to 0} \frac{f(x)}{g(x)}$ 不存在 $\underline{\hspace{3.5em}}$ （填“对”或“错”）．

  \vspace{1.2cm}
\end{problem}

% =========================================================================
% 三、计算解答题
% =========================================================================
\section*{三、计算解答题（解答应写出文字说明、演算步骤或推导过程。）}

\begin{problem}
  求函数 $f(x) = \frac{\ln|x|}{|x-1|} \sin x$ 的间断点并指出其类型．

  \vspace{6.0cm}
\end{problem}

\begin{problem}
  设 $\lim_{x \to \infty} (x^2+2x^4)^b - x^a = b$．且 $b \ne 0$，求常数 $a, b$．

  \vspace{6.0cm}
\end{problem}

% =========================================================================
% 参考答案与详细推导附录 (Solutions & Proofs Appendix)
% =========================================================================
\clearpage
\section*{参考答案与详细推导}

\subsection*{一、参考答案速查}

\begin{center}
\begin{tabular}{c p{5.5cm} c p{5.5cm}}
  \toprule
  \textbf{题号} & \textbf{参考答案} & \textbf{题号} & \textbf{参考答案} \\
  \midrule
  1 & D & 9 & A \\
  2 & A & 10 & C \\
  3 & D & 11 & $1$ \\
  4 & B & 12 & $\frac{6}{5}$ \\
  5 & D & 13 & 错 \\
  6 & C & 14 & 见详细解析 \\
  7 & B & 15 & $a = \frac{1}{2}$, $b = \frac{2}{3}$ \\
  8 & B &  &  \\
  \bottomrule
\end{tabular}
\end{center}

\subsection*{二、详细推导与证明过程}

\begin{solution}[第 11 题解答]
  \textbf{【答案】} $1$

\end{solution}

\begin{solution}[第 12 题解答]
  \textbf{【答案】} $\frac{6}{5}$

\end{solution}

\begin{solution}[第 13 题解答]
  \textbf{【答案】} 错

\end{solution}

\begin{solution}[第 14 题解答]
  \textbf{【答案】} $f(x)$ 有可去间断点 $x=0$，跳跃间断点 $x=1$

\end{solution}

\begin{solution}[第 15 题解答]
  \textbf{【答案】} $a = \frac{1}{2}$, $b = \frac{2}{3}$

\end{solution}

\end{document}
